\scanpage{514}
\mspage{426}
\deleted{
% B01
Alors \deleted{\unsure{les}} $i_s:E_s\unsure{\relbar|\longrightarrow}E$ satisfait \unsure{:|à} $L_{\pi_0}$

% B02
[En particulier, si $\pi_0$ est choisi de façon que $\pi_0$ majore tous les $\pi_\lambda$ ($\lambda\in\mathrm{Fl}\,C$), \deleted{alors} (c'est possible car $\mathrm{Fl}\,C$ est petit) alors $\forall s\in\mathrm{Ob}\,\unsure{S|C}$, $i_s:E_s\relbar E$ satisfait $L_{\pi_0}$.]

% B03
% margin note written sideways
\margin{$1)$ $\pi_0\geq\underset{\lambda\in\mathrm{Fl}\,C}{\mathrm{Sup}}\ \pi_\lambda$\\ $2)$ $\pi_1\geq\sup(\pi_0,\mathrm{card}\,C,\underset{s\in\mathrm{Ob}\,C}{\mathrm{Sup}}(\mathrm{card\,red}\,E_{s\pi_0})$.}

% B04
b) \inserted{Supposons $C$ \uline{sans projecteurs}.} \deleted{\unsure{Si} \unsure{tel} \unsure{quel} \unsure{que} \illegible{2} $\pi_0$ \unsure{sup} \illegible{1}} \inserted{$\pi_0\geq\underset{\lambda\in\mathrm{Fl}\,C}{\mathrm{Sup}}\pi_\lambda$,} \unsure{Soit|soient} $\pi_0,\pi_1$ tels que \deleted{\unsure{et soit} $\pi_1$ un cardinal $\unsure{\mathrm{card}\,C},\underset{s\in\mathrm{Ob}\,C}{\mathrm{Sup}}(\mathrm{card}\ \unsure{E_{s\pi_0}})$ \unsure{Sup} \illegible{2} $\mathrm{card}\,E_{s\pi_0}$} tel que \deleted{\illegible{2} $(\pi_0,\mathrm{card\,red}\,E_{s\pi_0})$, \unsure{satisfait au cardinal} $L_{\pi_1}$} Alors $E$ satisfait à $L_{\pi_1}$, et si $\pi$ est tel que \deleted{\illegible{1}} $\pi^{\pi_0}=\pi\geq\pi_1$, Alors
% B05
\[i_s(E_{s\pi})\subset E_\pi.\]

% B06
c) \inserted{\unsure{Sous les} hypothèses sur $C,\pi_0,\pi_1$ de \unsure{b)|b},} Pour tout cardinal $\pi$, soit $E(\pi)$ la sous-catégorie pleine de $E$ \deleted{\unsure{formée des}} telle que
% B07
\[\mathrm{Ob}\,E(\pi)=\bigcup_{s\in\mathrm{Ob}\,C}\mathrm{Ob}\,E_{s\pi}.\]

% B08
\deleted{Soit $\pi_0$ tel que $\pi_0\geq\underset{\lambda\in\mathrm{Fl}\,C}{\mathrm{Sup}}\pi_\lambda$, et $\pi_1\geq\pi_0$ un cardinal tel que $\pi_1\geq\underset{\deleted{\unsure{\lambda\in\mathrm{Fl}\,C|\mathrm{Fl}\,C}}\,s\in\mathrm{Ob}\,C}{\mathrm{Sup}}(\mathrm{card\,red}\,E_{s\pi_0})$, et que $E$ satisfasse $L_{\pi_1}$.}

% B09
Soit $\pi$ un cardinal tel que $\pi^{\pi_0}=\pi\geq\unsure{\pi_1|I_1}$, \deleted{et tel que $E$ satisfasse à $L_\pi$} Alors
% B10
\[E(\pi)\subset E_\pi,\]
% B11
d'où un foncteur pleinement fidèle ($E$ satisfaisant \unsure{:|à} $L_\pi$)
% B12
\[\mathrm{Ind}_\pi E(\pi)\longrightarrow E,\]
% B13
et ce foncteur est une équivalence (donc $E$ est accessible).

}
% archive sheet number: 513
